% Options for packages loaded elsewhere
\PassOptionsToPackage{unicode}{hyperref}
\PassOptionsToPackage{hyphens}{url}
\PassOptionsToPackage{dvipsnames,svgnames,x11names}{xcolor}
%
\documentclass[
  10pt,
  letterpaper,
]{article}

\usepackage{amsmath,amssymb}
\usepackage{iftex}
\ifPDFTeX
  \usepackage[T1]{fontenc}
  \usepackage[utf8]{inputenc}
  \usepackage{textcomp} % provide euro and other symbols
\else % if luatex or xetex
  \usepackage{unicode-math}
  \defaultfontfeatures{Scale=MatchLowercase}
  \defaultfontfeatures[\rmfamily]{Ligatures=TeX,Scale=1}
\fi
\usepackage{lmodern}
\ifPDFTeX\else  
    % xetex/luatex font selection
\fi
% Use upquote if available, for straight quotes in verbatim environments
\IfFileExists{upquote.sty}{\usepackage{upquote}}{}
\IfFileExists{microtype.sty}{% use microtype if available
  \usepackage[]{microtype}
  \UseMicrotypeSet[protrusion]{basicmath} % disable protrusion for tt fonts
}{}
\makeatletter
\@ifundefined{KOMAClassName}{% if non-KOMA class
  \IfFileExists{parskip.sty}{%
    \usepackage{parskip}
  }{% else
    \setlength{\parindent}{0pt}
    \setlength{\parskip}{6pt plus 2pt minus 1pt}}
}{% if KOMA class
  \KOMAoptions{parskip=half}}
\makeatother
\usepackage{xcolor}
\setlength{\emergencystretch}{3em} % prevent overfull lines
\setcounter{secnumdepth}{-\maxdimen} % remove section numbering
% Make \paragraph and \subparagraph free-standing
\ifx\paragraph\undefined\else
  \let\oldparagraph\paragraph
  \renewcommand{\paragraph}[1]{\oldparagraph{#1}\mbox{}}
\fi
\ifx\subparagraph\undefined\else
  \let\oldsubparagraph\subparagraph
  \renewcommand{\subparagraph}[1]{\oldsubparagraph{#1}\mbox{}}
\fi


\providecommand{\tightlist}{%
  \setlength{\itemsep}{0pt}\setlength{\parskip}{0pt}}\usepackage{longtable,booktabs,array}
\usepackage{calc} % for calculating minipage widths
% Correct order of tables after \paragraph or \subparagraph
\usepackage{etoolbox}
\makeatletter
\patchcmd\longtable{\par}{\if@noskipsec\mbox{}\fi\par}{}{}
\makeatother
% Allow footnotes in longtable head/foot
\IfFileExists{footnotehyper.sty}{\usepackage{footnotehyper}}{\usepackage{footnote}}
\makesavenoteenv{longtable}
\usepackage{graphicx}
\makeatletter
\def\maxwidth{\ifdim\Gin@nat@width>\linewidth\linewidth\else\Gin@nat@width\fi}
\def\maxheight{\ifdim\Gin@nat@height>\textheight\textheight\else\Gin@nat@height\fi}
\makeatother
% Scale images if necessary, so that they will not overflow the page
% margins by default, and it is still possible to overwrite the defaults
% using explicit options in \includegraphics[width, height, ...]{}
\setkeys{Gin}{width=\maxwidth,height=\maxheight,keepaspectratio}
% Set default figure placement to htbp
\makeatletter
\def\fps@figure{htbp}
\makeatother

\usepackage[preprint]{neurips_2025}
\usepackage{booktabs}
\usepackage{amsfonts}
\usepackage{amsmath}
\usepackage{nicefrac}
\usepackage{microtype}
\usepackage{xcolor}
\usepackage{graphicx}
\usepackage{subcaption}
\usepackage{multirow}
\usepackage{enumitem}
% Custom commands
\newcommand{\pcoop}{P(\text{C})}
\newcommand{\pcoopgiven}[1]{P(\text{C} \mid #1)}
\newcommand{\pp}{\text{\,pp}}
\newcommand{\fhat}{\hat{f}}
\newcommand{\commgap}{\Delta_{\text{comm}}}
\newcommand{\tablabel}[1]{\label{tab:#1}}
\setmainfont{texgyretermes-regular.otf}[
  Path=/usr/local/texlive/2024/texmf-dist/fonts/opentype/public/tex-gyre/,
  BoldFont=texgyretermes-bold.otf,
  ItalicFont=texgyretermes-italic.otf,
  BoldItalicFont=texgyretermes-bolditalic.otf
]
\setsansfont{texgyreheros-regular.otf}[
  Path=/usr/local/texlive/2024/texmf-dist/fonts/opentype/public/tex-gyre/,
  BoldFont=texgyreheros-bold.otf,
  ItalicFont=texgyreheros-italic.otf,
  BoldItalicFont=texgyreheros-bolditalic.otf
]
\setmonofont{texgyrecursor-regular.otf}[
  Path=/usr/local/texlive/2024/texmf-dist/fonts/opentype/public/tex-gyre/,
  BoldFont=texgyrecursor-bold.otf,
  ItalicFont=texgyrecursor-italic.otf,
  BoldItalicFont=texgyrecursor-bolditalic.otf
]
\setmathfont{texgyretermes-math.otf}[
  Path=/usr/local/texlive/2024/texmf-dist/fonts/opentype/public/tex-gyre-math/
]
\makeatletter
\makeatother
\makeatletter
\makeatother
\makeatletter
\@ifpackageloaded{caption}{}{\usepackage{caption}}
\AtBeginDocument{%
\ifdefined\contentsname
  \renewcommand*\contentsname{Table of contents}
\else
  \newcommand\contentsname{Table of contents}
\fi
\ifdefined\listfigurename
  \renewcommand*\listfigurename{List of Figures}
\else
  \newcommand\listfigurename{List of Figures}
\fi
\ifdefined\listtablename
  \renewcommand*\listtablename{List of Tables}
\else
  \newcommand\listtablename{List of Tables}
\fi
\ifdefined\figurename
  \renewcommand*\figurename{Figure}
\else
  \newcommand\figurename{Figure}
\fi
\ifdefined\tablename
  \renewcommand*\tablename{Table}
\else
  \newcommand\tablename{Table}
\fi
}
\@ifpackageloaded{float}{}{\usepackage{float}}
\floatstyle{ruled}
\@ifundefined{c@chapter}{\newfloat{codelisting}{h}{lop}}{\newfloat{codelisting}{h}{lop}[chapter]}
\floatname{codelisting}{Listing}
\newcommand*\listoflistings{\listof{codelisting}{List of Listings}}
\makeatother
\makeatletter
\@ifpackageloaded{caption}{}{\usepackage{caption}}
\@ifpackageloaded{subcaption}{}{\usepackage{subcaption}}
\makeatother
\makeatletter
\@ifpackageloaded{tcolorbox}{}{\usepackage[skins,breakable]{tcolorbox}}
\makeatother
\makeatletter
\@ifundefined{shadecolor}{\definecolor{shadecolor}{rgb}{.97, .97, .97}}
\makeatother
\makeatletter
\makeatother
\makeatletter
\makeatother
\ifLuaTeX
  \usepackage{selnolig}  % disable illegal ligatures
\fi
\usepackage[]{natbib}
\bibliographystyle{plainnat}
\IfFileExists{bookmark.sty}{\usepackage{bookmark}}{\usepackage{hyperref}}
\IfFileExists{xurl.sty}{\usepackage{xurl}}{} % add URL line breaks if available
\urlstyle{same} % disable monospaced font for URLs
\hypersetup{
  pdftitle={Same-Checkpoint Controls for Disentangling Communication Benefit from Learning Dynamics in Multi-Agent Reinforcement Learning},
  pdfauthor={Anonymous Authors},
  colorlinks=true,
  linkcolor={blue},
  filecolor={Maroon},
  citecolor={Blue},
  urlcolor={Blue},
  pdfcreator={LaTeX via pandoc}}

% Quarto template partial: NeurIPS title block
% Suppresses default Quarto author rendering; uses NeurIPS \maketitle
\title{Same-Checkpoint Controls for Disentangling Communication Benefit
from Learning Dynamics in Multi-Agent Reinforcement Learning}

\author{%
  Anonymous Authors \\
  % Author information omitted for double-blind review
}
\begin{document}
\maketitle
% Auto-generated by paper_data.py — do not edit manually
% Run: python paper_data.py --gen-defs

% Table 1: comm vs no-comm cooperation rates
\def\TIcommTFFk{0.775 \pm 0.028}  % comm f=3.500 50k
\def\TIncTFFk{0.472 \pm 0.012}  % no_comm f=3.500 50k
\def\TIgapTFFk{\mathbf{+0.303}}  % gap f=3.500 50k  p=0.0001
\def\TIcommTFHk{0.559 \pm 0.055}  % comm f=3.500 100k
\def\TIncTFHk{0.421 \pm 0.023}  % no_comm f=3.500 100k
\def\TIgapTFHk{+0.137}  % gap f=3.500 100k  p=0.0279
\def\TIcommTFOfk{0.536 \pm 0.038}  % comm f=3.500 150k
\def\TIncTFOfk{0.450 \pm 0.025}  % no_comm f=3.500 150k
\def\TIgapTFOfk{+0.086}  % gap f=3.500 150k  p=0.0609
\def\TIcommFVFk{0.761 \pm 0.032}  % comm f=5.000 50k
\def\TIncFVFk{0.732 \pm 0.038}  % no_comm f=5.000 50k
\def\TIgapFVFk{+0.029}  % gap f=5.000 50k  p=0.5427
\def\TIcommFVHk{0.556 \pm 0.064}  % comm f=5.000 100k
\def\TIncFVHk{0.649 \pm 0.042}  % no_comm f=5.000 100k
\def\TIgapFVHk{-0.093}  % gap f=5.000 100k  p=0.2766
\def\TIcommFVOfk{0.561 \pm 0.041}  % comm f=5.000 150k
\def\TIncFVOfk{0.800 \pm 0.032}  % no_comm f=5.000 150k
\def\TIgapFVOfk{-0.240}  % gap f=5.000 150k  p=0.0009

% Table 2: channel controls at 50k
\def\TIILearned{0.809}  % learned (n=5)
\def\TIIIndep{0.669}  % indep_random (n=3)
\def\TIIZero{0.641}  % always_zero (n=3)
\def\TIIPub{0.620}  % public_random (n=3)

% Table 3: continuation controls
\def\TIIInSeeds{15}  % number of seeds
\def\TIIIsource{same-checkpoint}  % data source type
\def\TIIIpending{}  % pending note
\def\TIIILearnedFk{0.775 \pm 0.028}  % learned 50k f=3.5
\def\TIIILearnedHk{0.559 \pm 0.055}  % learned 100k f=3.5
\def\TIIILearnedOfk{0.536 \pm 0.038}  % learned 150k f=3.5
\def\TIIILearnedFkFV{0.761 \pm 0.032}  % learned 50k f=5.0
\def\TIIILearnedOfkFV{0.561 \pm 0.041}  % learned 150k f=5.0
\def\TIIIIndepFk{0.775 \pm 0.028}  % indep_random 50k f=3.5
\def\TIIIIndepHk{0.459 \pm 0.028}  % indep_random 100k f=3.5
\def\TIIIIndepOfk{0.415 \pm 0.019}  % indep_random 150k f=3.5
\def\TIIIIndepFkFV{0.761 \pm 0.032}  % indep_random 50k f=5.0
\def\TIIIIndepOfkFV{0.631 \pm 0.044}  % indep_random 150k f=5.0
\def\TIIIIndepGapOfk{+12.1}  % learned - indep_random gap at 150k (pp)
\def\TIIIIndepGapOfkCI{[+3.4,\,+21.5]}  % 95% CI for gap (pp)
\def\TIIIIndepGapOfkP{=0.0189}  % p-value for gap
\def\TIIIPubFk{0.775 \pm 0.028}  % public_random 50k f=3.5
\def\TIIIPubHk{0.425 \pm 0.052}  % public_random 100k f=3.5
\def\TIIIPubOfk{0.399 \pm 0.022}  % public_random 150k f=3.5
\def\TIIIPubFkFV{0.761 \pm 0.032}  % public_random 50k f=5.0
\def\TIIIPubOfkFV{0.614 \pm 0.031}  % public_random 150k f=5.0
\def\TIIIPubGapOfk{+13.7}  % learned - public_random gap at 150k (pp)
\def\TIIIPubGapOfkCI{[+5.7,\,+21.5]}  % 95% CI for gap (pp)
\def\TIIIPubGapOfkP{=0.0072}  % p-value for gap
\def\TIIIZeroFk{0.775 \pm 0.028}  % always_zero 50k f=3.5
\def\TIIIZeroHk{0.343 \pm 0.048}  % always_zero 100k f=3.5
\def\TIIIZeroOfk{0.342 \pm 0.036}  % always_zero 150k f=3.5
\def\TIIIZeroFkFV{0.761 \pm 0.032}  % always_zero 50k f=5.0
\def\TIIIZeroOfkFV{0.544 \pm 0.043}  % always_zero 150k f=5.0
\def\TIIIZeroGapOfk{+19.4}  % learned - always_zero gap at 150k (pp)
\def\TIIIZeroGapOfkCI{[+11.3,\,+29.2]}  % 95% CI for gap (pp)
\def\TIIIZeroGapOfkP{=0.0002}  % p-value for gap

% Same-checkpoint sender-shuffle continuation (matched-stat control)
\def\TSSshufCoopTFOfk{0.446}  % sender-shuffle coop f=3.5 150k
\def\TSSshufCoopTFOfkSEM{0.028}  % sender-shuffle coop f=3.5 150k SEM
\def\TSSshufCoopTFOfkPM{0.446 \pm 0.028}  % sender-shuffle coop f=3.5 150k mean ± SEM
\def\TSSgapTFOfk{+9.0}  % learned - shuffle gap f=3.5 150k (pp)
\def\TSSgapTFOfkCI{[-1.0,\,+18.7]}  % 95% CI for gap f=3.5 150k (pp)
\def\TSSgapTFOfkSign{9/15}  % seeds favoring learned at f=3.5 150k
\def\TSSshufCoopFVOfk{0.720}  % sender-shuffle coop f=5.0 150k
\def\TSSshufCoopFVOfkSEM{0.038}  % sender-shuffle coop f=5.0 150k SEM
\def\TSSshufCoopFVOfkPM{0.720 \pm 0.038}  % sender-shuffle coop f=5.0 150k mean ± SEM
\def\TSSgapFVOfk{-15.9}  % learned - shuffle gap f=5.0 150k (pp)
\def\TSSgapFVOfkCI{[-29.4,\,-3.4]}  % 95% CI for gap f=5.0 150k (pp)
\def\TSSgapFVOfkSign{6/15}  % seeds favoring learned at f=5.0 150k
\def\TIIIminOfk{0.342}  % min 150k across modes
\def\TIIImaxOfk{0.536}  % max 150k across modes

% Table 4: frozen-checkpoint interventions
\def\TFRZerosFk{+0.0}  % zeros 50k (pp)
\def\TFRZerosOfk{+6.1}  % zeros 150k (pp)
\def\TFRZerosFkSEM{2.4}  % zeros 50k SEM (pp)
\def\TFRZerosOfkSEM{4.2}  % zeros 150k SEM (pp)
\def\TFRZerosOfkCI{[-1.2,\,+14.5]}  % zeros 150k 95% CI (pp)
\def\TFRZerosOfkP{=0.1667}  % zeros 150k p-value
\def\TFRZerosOfkSign{9/15}  % zeros 150k seeds with natural > intervention
\def\TFRFixZeroFk{-2.2}  % fixed0 50k (pp)
\def\TFRFixZeroOfk{+3.0}  % fixed0 150k (pp)
\def\TFRFixZeroFkSEM{3.2}  % fixed0 50k SEM (pp)
\def\TFRFixZeroOfkSEM{3.2}  % fixed0 150k SEM (pp)
\def\TFRFixZeroOfkCI{[-3.1,\,+8.8]}  % fixed0 150k 95% CI (pp)
\def\TFRFixZeroOfkP{=0.3494}  % fixed0 150k p-value
\def\TFRFixZeroOfkSign{9/15}  % fixed0 150k seeds with natural > intervention
\def\TFRPubFk{-2.0}  % public_random 50k (pp)
\def\TFRPubOfk{+1.6}  % public_random 150k (pp)
\def\TFRPubFkSEM{1.5}  % public_random 50k SEM (pp)
\def\TFRPubOfkSEM{1.6}  % public_random 150k SEM (pp)
\def\TFRPubOfkCI{[-1.3,\,+4.8]}  % public_random 150k 95% CI (pp)
\def\TFRPubOfkP{=0.3741}  % public_random 150k p-value
\def\TFRPubOfkSign{7/15}  % public_random 150k seeds with natural > intervention
\def\TFRIndepFk{-0.3}  % indep_random 50k (pp)
\def\TFRIndepOfk{+1.6}  % indep_random 150k (pp)
\def\TFRIndepFkSEM{0.7}  % indep_random 50k SEM (pp)
\def\TFRIndepOfkSEM{0.7}  % indep_random 150k SEM (pp)
\def\TFRIndepOfkCI{[+0.4,\,+3.0]}  % indep_random 150k 95% CI (pp)
\def\TFRIndepOfkP{=0.0286}  % indep_random 150k p-value
\def\TFRIndepOfkSign{10/15}  % indep_random 150k seeds with natural > intervention
\def\TFRShuffleFk{-2.9}  % sender_shuffle 50k (pp)
\def\TFRShuffleOfk{-3.2}  % sender_shuffle 150k (pp)
\def\TFRShuffleFkSEM{0.9}  % sender_shuffle 50k SEM (pp)
\def\TFRShuffleOfkSEM{0.6}  % sender_shuffle 150k SEM (pp)
\def\TFRShuffleOfkCI{[-4.4,\,-2.3]}  % sender_shuffle 150k 95% CI (pp)
\def\TFRShuffleOfkP{<0.0001}  % sender_shuffle 150k p-value
\def\TFRShuffleOfkSign{0/15}  % sender_shuffle 150k seeds with natural > intervention
\def\TFRPermuteFk{-0.4}  % permute_slots 50k (pp)
\def\TFRPermuteOfk{+1.6}  % permute_slots 150k (pp)
\def\TFRPermuteFkSEM{0.5}  % permute_slots 50k SEM (pp)
\def\TFRPermuteOfkSEM{0.7}  % permute_slots 150k SEM (pp)
\def\TFRPermuteOfkCI{[+0.3,\,+3.1]}  % permute_slots 150k 95% CI (pp)
\def\TFRPermuteOfkP{=0.0461}  % permute_slots 150k p-value
\def\TFRPermuteOfkSign{11/15}  % permute_slots 150k seeds with natural > intervention

% Table 5: endpoint sender-causal probe
\def\TSCDeltaZeroFive{-0.0}  % sender-causal mean delta at f=0.5 (pp)
\def\TSCAbsZeroFive{0.1}  % sender-causal mean abs delta at f=0.5 (pp)
\def\TSCFlipZeroZeroFive{0.1}  % sender-causal flip0 at f=0.5 (%)
\def\TSCFlipOneZeroFive{0.1}  % sender-causal flip1 at f=0.5 (%)
\def\TSCDeltaOneFive{-0.0}  % sender-causal mean delta at f=1.5 (pp)
\def\TSCAbsOneFive{0.9}  % sender-causal mean abs delta at f=1.5 (pp)
\def\TSCFlipZeroOneFive{0.7}  % sender-causal flip0 at f=1.5 (%)
\def\TSCFlipOneOneFive{0.6}  % sender-causal flip1 at f=1.5 (%)
\def\TSCDeltaTwoFive{-0.7}  % sender-causal mean delta at f=2.5 (pp)
\def\TSCAbsTwoFive{7.7}  % sender-causal mean abs delta at f=2.5 (pp)
\def\TSCFlipZeroTwoFive{7.3}  % sender-causal flip0 at f=2.5 (%)
\def\TSCFlipOneTwoFive{6.9}  % sender-causal flip1 at f=2.5 (%)
\def\TSCDeltaThreeFive{+0.3}  % sender-causal mean delta at f=3.5 (pp)
\def\TSCAbsThreeFive{5.0}  % sender-causal mean abs delta at f=3.5 (pp)
\def\TSCFlipZeroThreeFive{12.2}  % sender-causal flip0 at f=3.5 (%)
\def\TSCFlipOneThreeFive{12.9}  % sender-causal flip1 at f=3.5 (%)
\def\TSCDeltaFiveZero{+0.4}  % sender-causal mean delta at f=5.0 (pp)
\def\TSCAbsFiveZero{3.8}  % sender-causal mean abs delta at f=5.0 (pp)
\def\TSCFlipZeroFiveZero{9.8}  % sender-causal flip0 at f=5.0 (%)
\def\TSCFlipOneFiveZero{11.5}  % sender-causal flip1 at f=5.0 (%)

% Table 6: mute experiment
\def\TIVcommCoop{0.536}  % comm coop at 150k
\def\TIVcommWelf{37.43}  % comm welfare
\def\TIVncCoop{0.450}  % no-comm coop at 150k
\def\TIVncWelf{34.01}  % no-comm welfare
\def\TIVmuteCoop{0.367}  % mute coop at 150k
\def\TIVmuteWelf{30.69}  % mute welfare
\def\TIVmuteCoopPct{36.7}  % mute coop %
\def\TIVncCoopPct{45.0}  % no-comm coop %

% Table 7: aggregate vs per-sender token effects
\def\TVaggFk{-0.000}  % aggregate 50k
\def\TVpsFk{0.218}  % per-sender 50k
\def\TVratioFk{---}  % ratio 50k (agg~0)
\def\TVaggHk{-0.008}  % aggregate 100k
\def\TVpsHk{0.178}  % per-sender 100k
\def\TVratioHk{23.3\times}  % ratio 100k
\def\TVaggOfk{0.003}  % aggregate 150k
\def\TVpsOfk{0.186}  % per-sender 150k
\def\TVratioOfk{---}  % ratio 150k (agg~0)

% Table 8: polarity alignment (summary over 15 seeds)
\def\TVInSeeds{15}  % number of seeds in alignment table
\def\TVItrajAchievedfull{2}  % n seeds with trajectory 'Achieved full alignment'
\def\TVItrajImproved{3}  % n seeds with trajectory 'Improved'
\def\TVItrajLostalignmen{5}  % n seeds with trajectory 'Lost alignment'
\def\TVItrajPolarityflip{2}  % n seeds with trajectory 'Polarity flipped'
\def\TVItrajStable{3}  % n seeds with trajectory 'Stable'

% Prose claims
\def\pvGapTFFkPP{30.3}  % comm gap f=3.5 50k (pp)
\def\pvGapTFOfkPP{8.6}  % comm gap f=3.5 150k (pp)
\def\pvGapTFFkPval{=0.0001}  % p-value for gap at f=3.5 50k
\def\pvGapDeclinePct{72}  % gap decline %
\def\pvGapRangeLo{10.5}  % per-seed gap min at 50k (pp)
\def\pvGapRangeHi{49.4}  % per-seed gap max at 50k (pp)
\def\pvNposOfk{10}  % n seeds with positive gap at 150k
\def\pvGapFVFkPP{2.9}  % comm gap f=5.0 50k (pp)
\def\pvGapFVOfkPP{-24.0}  % comm gap f=5.0 150k (pp)
\def\pvCtrlAdvLo{14}  % control advantage low (pp)
\def\pvCtrlAdvHi{19}  % control advantage high (pp)
\def\pvMuteBelowPP{8}  % mute below no-comm (pp)
\def\pvRatioFk{>44}  % per-sender/agg ratio 50k
\def\pvRatioOfk{---}  % per-sender/agg ratio 150k
\def\pvAggFkPP{-0.0}  % aggregate token effect 50k (pp)
\def\pvPsFkPP{21.8}  % per-sender effect 50k (pp)
\def\pvAggOfkPP{0.3}  % aggregate effect 150k (pp)
\def\pvPsOfkPP{18.6}  % per-sender effect 150k (pp)
\def\pvCommCoopTFFkPct{77.5}  % comm coop % at f=3.5 50k
\def\pvNcCoopTFFkPct{47}  % no-comm coop % at f=3.5 50k
\def\pvNSeeds{15}  % total number of seeds

% Appendix: summary stats for cooperation and gaps at f=3.5
\def\pvPSnSeedsFk{15}  % n seeds at 50k
\def\pvPScommMeanFk{0.775}  % mean comm 3.5 50k
\def\pvPScommSDFk{0.108}  % SD comm 3.5 50k
\def\pvPScommMinFk{0.544}  % min comm 3.5 50k
\def\pvPScommMaxFk{0.904}  % max comm 3.5 50k
\def\pvPSgapMeanFk{+0.303}  % mean gap 3.5 50k
\def\pvPSgapSDFk{0.094}  % SD gap 3.5 50k
\def\pvPSgapMinFk{+0.105}  % min gap 3.5 50k
\def\pvPSgapMaxFk{+0.494}  % max gap 3.5 50k
\def\pvPSnSeedsHk{15}  % n seeds at 100k
\def\pvPScommMeanHk{0.559}  % mean comm 3.5 100k
\def\pvPScommSDHk{0.211}  % SD comm 3.5 100k
\def\pvPScommMinHk{0.219}  % min comm 3.5 100k
\def\pvPScommMaxHk{0.919}  % max comm 3.5 100k
\def\pvPSgapMeanHk{+0.137}  % mean gap 3.5 100k
\def\pvPSgapSDHk{0.218}  % SD gap 3.5 100k
\def\pvPSgapMinHk{-0.233}  % min gap 3.5 100k
\def\pvPSgapMaxHk{+0.544}  % max gap 3.5 100k
\def\pvPSnSeedsOfk{15}  % n seeds at 150k
\def\pvPScommMeanOfk{0.536}  % mean comm 3.5 150k
\def\pvPScommSDOfk{0.146}  % SD comm 3.5 150k
\def\pvPScommMinOfk{0.325}  % min comm 3.5 150k
\def\pvPScommMaxOfk{0.890}  % max comm 3.5 150k
\def\pvPSgapMeanOfk{+0.086}  % mean gap 3.5 150k
\def\pvPSgapSDOfk{0.162}  % SD gap 3.5 150k
\def\pvPSgapMinOfk{-0.167}  % min gap 3.5 150k
\def\pvPSgapMaxOfk{+0.363}  % max gap 3.5 150k


\ifdefined\Shaded\renewenvironment{Shaded}{\begin{tcolorbox}[breakable, borderline west={3pt}{0pt}{shadecolor}, frame hidden, interior hidden, enhanced, boxrule=0pt, sharp corners]}{\end{tcolorbox}}\fi

\begin{abstract}
Emergent communication in multi-agent reinforcement learning is often evaluated with endpoint metrics such as cooperation, mutual information, or message--action correlations measured at the final checkpoint. Under independent decentralized learning, however, such metrics can conflate two distinct questions: whether message content shaped the training trajectory, and whether the frozen endpoint policy uses the online message stream as a robust, sender-consistent, and behaviorally useful code. We introduce \emph{same-checkpoint continuation controls} to separate these questions: identical learned checkpoints are forked into learned, random, and silent channel modes, and their downstream outcomes are compared during continued training. As a case study, we apply this design to a four-agent extended public goods game with hidden regime switching and noisy private observations (15 seeds). Same-checkpoint controls show that learned messages retain a clear continued-learning advantage in the mixed-motive regime, outperforming silence and random channels by roughly 12--19 percentage points, while all channel modes perform similarly in the cooperation-dominant regime. A matched-stat sender-shuffle continuation control shrinks but does not erase the mixed-regime advantage, suggesting a mixture of genuine state-linked message value and broader channel dependence. Complementary frozen-checkpoint interventions show that this training-time benefit does not translate into a robust endpoint protocol: sender-shuffled messages outperform natural messages, and sender-specific causal probes reveal sizable local effects that cancel in aggregate. These results show that communication benefit and learning dynamics can be more cleanly disentangled, and that continued-learning value and endpoint protocol coherence can come apart in ways that endpoint-only evaluation misses.
\end{abstract}

\hypertarget{sec-intro}{%
\section{Introduction}\label{sec-intro}}

Multi-agent reinforcement learning (MARL) has produced a growing body of
work on emergent communication, in which agents learn to use a discrete
signaling channel to coordinate behavior
\citep{foerster2016learning, sukhbaatar2016learning, lazaridou2017multi, havrylov2017emergence, mordatch2018emergence}.
For a comprehensive survey, see \citet{lazaridou2020emergent}. Beyond
classic sender--receiver benchmarks, communication is now studied in
broader cooperative settings, including sequential social dilemmas,
intention-sharing architectures, and zero-shot convention formation
\citep{leibo2017multi, kim2021communication, hu2020other, bullard2021quasi}.
A central question is whether agents develop meaningful, shared
communication protocols or merely learn policies that happen to
correlate with message-sending behavior
\citep{kottur2017natural, cao2018emergent}. Standard evaluation practice
therefore focuses on \emph{endpoint metrics}: mutual information (MI)
between messages and task-relevant variables \citep{lowe2019pitfalls},
topographic similarity \citep{lazaridou2017multi}, or cooperation rates
measured at the final checkpoint.

This attribution problem is especially acute under independent learning,
where non-stationarity makes it difficult to separate message semantics
from the broader training dynamics
\citep{papoudakis2019dealing, oroojlooyjadid2022review}. We study it in
a four-agent extended public goods game (EPGG) with hidden regime
switching, noisy private observations, and optional binary communication
\citep{orzan2024pgg}. The EPGG is a case study for this measurement
problem rather than the contribution by itself: it gives a setting with
fixed teams, sender-indexed message slots, feed-forward policies, and
mixed incentives in which endpoint-only evaluation is especially likely
to mislead.

The setting sits at the intersection of several adjacent literatures.
Partially aligned signaling with hidden state is the core concern in
classic cheap-talk and strategic information transmission models
\citep{crawford1982strategic, chen2008selecting}. Communication in
public goods games also connects to experimental and institutional work
on cooperation under uncertainty
\citep{ledyard1995public, fehr2000cooperation, kosfeld2009institution, orzan2024pgg}.
More broadly, the paper fits within the Cooperative AI agenda of
understanding how decentralized agents and institutions can achieve
robust cooperation under imperfect alignment and limited information
\citep{dafoe2020open}.

We intentionally begin with independent learning rather than centralized
training. The goal is not to maximize coordination with the strongest
available training machinery, but to study a regime where attribution is
genuinely difficult: each agent learns under non-stationarity from local
observations and the message channel itself. In this setting,
improvements can easily be misread as endpoint semantics when they may
instead reflect broader co-learning dynamics. Independent PPO is
therefore a demanding but appropriate starting point for the evaluation
question we ask here.

We argue that endpoint evaluation alone can obscure continued-learning
value and overstate endpoint protocol coherence. Throughout, by
\emph{endpoint protocol coherence} we mean whether the frozen endpoint
policy satisfies three properties: (i)\textasciitilde it benefits from
the natural message stream relative to perturbations
(\emph{robustness}), (ii)\textasciitilde receivers respond consistently
to a given token regardless of which sender produced it
(\emph{sender-consistency}), and (iii)\textasciitilde the net effect of
natural messages on cooperation is positive (\emph{behavioral benefit}).
A protocol that is locally coherent---each sender--receiver pair has a
consistent convention---but globally fragmented across senders would
fail properties (ii) and potentially (iii). More generally, three
failure modes are worth distinguishing:

\begin{enumerate}[leftmargin=*,nosep]
  \item \textbf{Aggregate cancellation.} When agents develop sender-indexed conventions with opposing polarities (e.g., agent~$i$ sends ``1'' for high regimes, agent~$j$ sends ``1'' for low regimes), aggregate metrics average to near zero even though each agent--receiver pair maintains a strong sender-conditioned relationship.
  \item \textbf{Strategic value-of-information gating.} Communication may matter most when incentives are mixed and the latent state is ambiguous, but add little where cooperation is already individually optimal, so evaluation at a single difficulty level or averaged across regimes can obscure when messages actually matter.
  \item \textbf{Structural dependence without internalization.} The trained policy may depend on the \emph{presence} of a message channel (as an architectural input) without having internalized cooperative strategies---removing the channel produces worse outcomes than never having communicated.
\end{enumerate}

The EPGG is useful here because the severity of the free-rider problem
changes over time without direct notification. Communication is
therefore potentially valuable, but not equally so in every regime.

Our contributions are:

\begin{enumerate}[leftmargin=*,nosep]
  \item \textbf{Same-checkpoint continuation controls.} We introduce a clean experimental design for isolating the role of message content during continued learning: fork identical learned 50k checkpoints into learned, random, or silent channel modes and compare downstream cooperation. This separates continued-learning channel-content effects from frozen-endpoint protocol behavior within a fixed architecture and task.
  \item \textbf{Continued-learning benefit vs.\ weak endpoint coherence.} Applying this design to a four-agent EPGG with 15~seeds, we find that learned messages retain a clear advantage ($+12$--$19\pp$) in the mixed-motive regime during continued training from 50k to 150k, while frozen-checkpoint interventions reveal a much weaker endpoint protocol: sender-shuffled messages \emph{outperform} natural messages across all 15~seeds.
  \item \textbf{Strategic value-of-information in upper mixed regimes.} In this EPGG setting, communication matters where incentive alignment is ambiguous ($f\!=\!2.5, 3.5$), but the same-checkpoint controls show no clean channel-content effect where cooperation is already individually optimal ($f\!=\!5.0$).
  \item \textbf{Endpoint causal disaggregation.} Frozen sender-specific token interventions show that endpoint messages are not inert: signed aggregate effects remain near zero while mean absolute sender--receiver effects remain sizeable, revealing sender-indexed conventions that are locally causal but globally incompatible.
\end{enumerate}

\hypertarget{sec-setup}{%
\section{Setup}\label{sec-setup}}

\hypertarget{extended-public-goods-game-with-regime-switching}{%
\subsection{Extended public goods game with regime
switching}\label{extended-public-goods-game-with-regime-switching}}

We study a four-agent (\(N=4\)) Extended Public Goods Game (EPGG) played
over \(T=100\) rounds per episode. Each round \(t\), agent \(i\)
observes a noisy private signal
\(\fhat_{i,t} = f_t + \varepsilon_{i,t}\), where
\(\varepsilon_{i,t} \sim \mathcal{N}(0, \sigma^2)\) with
\(\sigma = 0.5\), and chooses to cooperate (\(a_i = C\)) or defect
(\(a_i = D\)). Each agent holds an endowment of \(c=4\) units. A
cooperator contributes \(c\) to a shared pool; a defector keeps \(c\) as
private payoff. Contributions are multiplied by the current regime
factor \(f_t\) and split equally, so each agent's round payoff is:
\begin{equation}\protect\hypertarget{eq-payoff}{}{
  u_i = \frac{f_t}{N} \sum_{j=1}^{N} c \cdot \mathbf{1}[a_j = C] \;+\; c \cdot \mathbf{1}[a_i = D]
}\label{eq-payoff}\end{equation} where the first term is agent \(i\)'s
share of the public pool and the second is the private endowment
retained when defecting (zero when cooperating). Throughout, we report
effect sizes in percentage points (pp): a change from 0.40 to 0.53 is
\(+13\pp\).

The hidden multiplier \(f_t\) takes values in
\(\mathcal{F} = \{0.5, 1.5, 2.5, 3.5, 5.0\}\) and switches with
probability \(\rho = 0.05\) per round (geometric holding time
\({\sim}20\) rounds). Each agent's action is subject to trembling: with
probability \(\varepsilon_{\text{tremble}} = 0.05\), the action is
replaced by a uniform random choice, modeling execution noise.

When \(f < N = 4\), defection is the dominant strategy; when \(f > N\),
cooperation dominates regardless of others' actions. The mixed-motive
regime (\(f = 3.5\)) is the critical case: cooperation is collectively
beneficial but individually costly (\(f/N = 0.875 < 1\), so each unit
contributed returns less than one unit to the contributor). This regime
creates the strongest demand for communication. The full multiplier set
spans five regimes, but the contrast between \(f\!=\!3.5\) and
\(f\!=\!5.0\) is theoretically central because it separates ambiguous
mixed incentives from a regime where cooperation is individually
optimal. Unless otherwise noted, reported quantities such as
\(\pcoopgiven{f}\) condition on the \emph{true round-level multiplier}
\(f_t\) during evaluation rather than on an episode-level label.

\hypertarget{communication-channel}{%
\subsection{Communication channel}\label{communication-channel}}

In the communication condition (\texttt{comm\_symm}), agents
simultaneously broadcast a binary message \(m_i \in \{0, 1\}\) before
choosing actions. All agents observe the same sender-indexed message
vector. Messages have no pre-assigned meaning; any semantics must emerge
through learning. In the no-communication condition
(\texttt{no\_comm\_symm}), no messages are exchanged. Agents never
observe the true multiplier \(f_t\), other agents' private signals, or
other agents' rewards.

\hypertarget{training}{%
\subsection{Training}\label{training}}

We train agents using independent PPO \citep{schulman2017proximal} with
generalized advantage estimation \citep{schulman2016high}
(\(\gamma = 0.99\), \(\lambda = 0.95\), clip ratio \(= 0.2\)). Each
agent has separate two-layer MLPs for the action policy, message policy
(for senders), and value prediction, each with hidden width 64. At each
round, a sender first emits a message from its current observation, and
all agents then select actions conditioned on the received same-round
message vector. For sender-updated samples, the action head and message
head are optimized jointly with one clipped objective over
\(\log \pi(a_t, m_t)\), using the agent's own reward return and the same
GAE advantage; the value head predicts only that agent's own discounted
return. The policies are not recurrent. Instead, temporal information
enters through a trainer-side observation wrapper that converts the raw
environment observation \([\fhat, \text{endowment}]\) into Set-A inputs
\([\fhat, \text{endowment}, \text{last-round cooperation fraction}, \text{own last action}, \text{EWMA cooperation}]\),
where EWMA denotes an exponential moving average, plus one two-entry
one-hot block per sender when communication is enabled. The base
observation is therefore 5-dimensional; with four senders and binary
messages, the communication slice adds eight dimensions (two per
sender), giving a total observation size of 13. Agents therefore track
the sticky latent regime only through current noisy signals, these
simple lagged summaries, and same-round communication. Key design
choices:

\textbf{Why independent PPO.} This choice is deliberate rather than
accidental. A centralized critic or stronger shared-training signal
could partly solve the coordination problem during training, making it
harder to tell whether later gains reflect decentralized online message
use or richer training-time credit assignment. We therefore start from
the harder decentralized regime in which endpoint-only evaluation is
most likely to be misleading.

\textbf{Entropy annealing.} Both action and message entropy coefficients
decay linearly over training: action entropy from \(0.01 \to 0.001\),
message entropy from \(0.01 \to 0.0\). This allows early exploration
while permitting conventions to crystallize. The learning rate follows a
cosine schedule from \(3 \times 10^{-4}\) to \(10^{-5}\).

\textbf{Seeds.} Fifteen independent seeds per condition, trained for
150k episodes with checkpoints at 50k, 100k, and 150k.

\hypertarget{sec-eval}{%
\subsection{Evaluation suite}\label{sec-eval}}

All evaluation uses greedy action selection (argmax policy) to isolate
learned behavior from exploration noise. These greedy checkpoint
evaluations are our main evidence about policy quality over training. By
contrast, the dense online training curves introduced later are
descriptive context only: they summarize sampled on-policy behavior
during training and therefore still reflect tremble noise, message
dropout, and residual exploration. Unless otherwise noted, means, SEMs,
bootstrap confidence intervals, and sign-flip tests treat training seed
as the unit of analysis; paired tests operate on per-seed paired deltas.
We evaluate four diagnostic dimensions:

\textbf{Message interventions.} At each checkpoint, we freeze the policy
and replay episodes while tampering with the messages agents receive.
The interventions test different hypotheses about what matters in the
message stream:

\begin{itemize}[leftmargin=*,nosep]
  \item \texttt{zeros}: zero the entire message slice of the wrapped observation after message construction---tests whether the policy benefits from \emph{any} message input.
  \item \texttt{fixed0}: replace every sender slot with token~0 before one-hot encoding---tests a constant but non-empty message pattern.
  \item \texttt{indep\_random}: replace sender slots with independent random bits each round---tests whether slot-wise noise is enough.
  \item \texttt{public\_random}: replace all sender slots with the same random bit each round---tests whether a shared coordination cue suffices.
  \item \texttt{sender\_shuffle}: for each sender, replace the current token with one drawn from that sender's empirical token pool---preserving per-sender marginal token frequency while breaking alignment with current state and time.
  \item \texttt{permute\_slots}: randomly permute sender slots each round while preserving the delivered bits---tests whether receivers track \emph{who} sent a message.
\end{itemize}

These frozen-checkpoint perturbations estimate the effect of specific
endpoint message manipulations; they do not, by themselves, identify the
full causal value of communication over training.

\textbf{Channel controls.} We compare four channel modes during
\emph{continued training} from the same learned 50k checkpoint:
\texttt{learned} (sender's actual message), \texttt{always\_zero} (all
sender slots forced to token\textasciitilde0), \texttt{indep\_random}
(sender slots randomized independently each round), and
\texttt{public\_random} (all sender slots share the same random bit each
round). This tests whether message \emph{content} matters for continued
learning once initialization is held fixed.

\textbf{Sender and receiver semantics.} To summarize how message use
varies with the perceived regime, we discretize each agent's noisy
private observation \(\fhat\) into five bins, denoted
\(\fhat_{\text{bin}}\): \(\fhat < 1.5\) (low), \(1.5 \le \fhat < 2.5\),
\(2.5 \le \fhat < 3.5\), \(3.5 \le \fhat < 4.5\), and \(\fhat \ge 4.5\)
(high). We then compute \(P(m_i = 1 \mid \fhat_{\text{bin}})\) (sender
semantics: how often sender \(i\) sends token\textasciitilde1 in each
bin) and \(P(a_j = C \mid m_i, \fhat_{\text{bin}})\) (receiver
semantics: how receiver \(j\)'s cooperation varies with sender \(i\)'s
token), both disaggregated by sender identity. When we later refer to
low and high bins for polarity, we mean \(\fhat < 1.5\) and
\(\fhat \ge 4.5\), respectively. This binned notation
\(\fhat_{\text{bin}}\) reappears in Appendix\textasciitilde C.

\textbf{Sender-specific causal endpoint probe.} At the frozen 150k
endpoint, we take greedy-evaluation states from the natural rollout,
hold the raw observations and all other message slots fixed, force
sender \(i\)'s token to 0 or 1, and measure the induced change in
receiver \(j\)'s cooperation probability and greedy action. Unlike the
receiver-semantics summaries, this is an interventional test of whether
sender-conditioned endpoint effects are causally used by the frozen
policy.

These diagnostics answer different questions. Same-checkpoint
continuation isolates whether channel content affects
\emph{continued learning} when initialization is fixed. Frozen message
interventions ask whether the \emph{endpoint policy} benefits from the
natural online message stream. Sender-specific frozen interventions ask
whether individual sender tokens retain local causal influence at the
endpoint. Receiver-semantics analyses are descriptive rather than causal
and are included to show how aggregate endpoint summaries can hide
sender-conditioned structure.

\textbf{Roadmap.} The remaining sections are organized as follows. The
\protect\hyperlink{sec-main-result}{baseline comparison} establishes a
descriptive baseline by comparing separately trained communication and
no-communication policies. The
\protect\hyperlink{sec-semantic-window}{continued-learning analysis}
presents the paper's core identification design:
\protect\hyperlink{sec-sameckpt}{same-checkpoint continuation} and
\protect\hyperlink{sec-frozen-endpoint}{frozen-checkpoint endpoint
interventions}, with exploratory controls deferred to
Appendices\textasciitilde D--E. The
\protect\hyperlink{sec-discussion}{Discussion} then places the results
in context and discusses limitations, while Appendix\textasciitilde C
provides the observational aggregate-metric and polarity analysis that
complements the frozen probes.

\hypertarget{sec-main-result}{%
\section{Baseline communication gap in separately trained
policies}\label{sec-main-result}}

\begin{table}[t]
  \caption{Round-conditional cooperation rates $P(C \mid f_t=f)$ for communication vs.\ no-communication conditions across training. Mean $\pm$ SEM over 15 seeds. The final row reports the paired mean gap (\texttt{comm\_symm} minus \texttt{no\_comm\_symm}).}
  \tablabel{comm_gap}
  \centering
  \small
  \begin{tabular}{llccc}
    \toprule
    & & \multicolumn{3}{c}{Checkpoint (episodes)} \\
    \cmidrule(lr){3-5}
    $f$ regime & Condition & 50k & 100k & 150k \\
    \midrule
    \multirow{3}{*}{$f\!=\!3.5$}
      & \texttt{comm\_symm}     & $\TIcommTFFk$ & $\TIcommTFHk$ & $\TIcommTFOfk$ \\
      & \texttt{no\_comm\_symm} & $\TIncTFFk$ & $\TIncTFHk$ & $\TIncTFOfk$ \\
      & $\commgap$              & $\TIgapTFFk$ & $\TIgapTFHk$ & $\TIgapTFOfk$ \\
    \midrule
    \multirow{3}{*}{$f\!=\!5.0$}
      & \texttt{comm\_symm}     & $\TIcommFVFk$ & $\TIcommFVHk$ & $\TIcommFVOfk$ \\
      & \texttt{no\_comm\_symm} & $\TIncFVFk$ & $\TIncFVHk$ & $\TIncFVOfk$ \\
      & $\commgap$              & $\TIgapFVFk$ & $\TIgapFVHk$ & $\TIgapFVOfk$ \\
    \bottomrule
  \end{tabular}
\end{table}

Table \ref{tab:comm_gap} is a descriptive baseline rather than the
paper's main identification design: communication and no-communication
policies are trained separately, so the gap mixes channel-content
effects with broader training-trajectory differences. We use it to
motivate the tighter same-checkpoint controls below, and
Figure~\ref{fig-value-surface} later shows that the descriptive gap
peaks in the upper mixed regimes rather than increasing monotonically
with generic ``difficulty.'\,'

At \(f = 3.5\) (mixed regime), communication produces a robust
\(+\pvGapTFFkPP\pp\) cooperation advantage at 50k episodes, with all 15
seeds showing positive effects (range: \(+\pvGapRangeLo\) to
\(+\pvGapRangeHi\pp\)). This benefit persists more modestly at 150k
(\(+\pvGapTFOfkPP\pp\), \(\pvNposOfk\)/15 seeds positive).

At \(f = 5.0\) (cooperation-dominant regime), the communication benefit
is small at 50k (\(+\pvGapFVFkPP\pp\)) and \textbf{reverses} by 150k
(\(\pvGapFVOfkPP\pp\), \(p < 0.001\)). In this EPGG setting,
communication-trained agents cooperate less than those that never had a
channel when cooperation is individually rational. This reversal should
be interpreted cautiously: Table \ref{tab:comm_gap} compares separately
trained communication and no-communication pipelines, so the gap could
reflect different training trajectories or architectural dependence on
the channel rather than harmful message content per se. Two mechanisms
could explain why having a channel during training \emph{hurts} at
\(f\!=\!5.0\): (i)\textasciitilde the communication architecture adds
input dimensions and policy complexity that slow convergence toward the
individually rational cooperation strategy, and
(ii)\textasciitilde mixed-motive regimes (\(f\!=\!3.5\)) create gradient
pressure toward conditional cooperation, and this conditioning may not
fully disengage when the regime switches to \(f\!=\!5.0\). The
same-checkpoint controls show no channel-mode separation at
\(f\!=\!5.0\) (Table \ref{tab:sameckpt}), confirming that the reversal
reflects these training-trajectory differences rather than ongoing
harmful message-content effects.

Welfare moves in the same direction, so the descriptive gap is not
driven by cooperation alone. Communication helps most where latent-state
information matters under mixed incentives, and the \(f\!=\!3.5\)
benefit declines from 50k to 150k without disappearing entirely.
Appendix\textasciitilde F shows dense online curves with checkpoint
anchors; we treat those trajectories as descriptive context only because
they cross a continuation seam and mix sampled training behavior with
greedy checkpoint evaluation.

\hypertarget{sec-semantic-window}{%
\section{When does message content matter?}\label{sec-semantic-window}}

The communication benefit raises a question: does message \emph{content}
drive the cooperation advantage, or does merely having a channel (even a
noisy one) suffice? We address this with a sequence of controls centered
on same-checkpoint continuation.

\hypertarget{sec-separate-controls}{%
\subsection{Exploratory from-scratch channel
controls}\label{sec-separate-controls}}

\hypertarget{sec-sameckpt}{%
\subsection{Same-checkpoint continuation isolates continued-training
channel-content effects}\label{sec-sameckpt}}

Appendix\textasciitilde D reports a small-(N) from-scratch
channel-control check, but because it confounds channel mode with the
entire training trajectory, we do not rely on it for identification.
Instead, we take the \textbf{identical} learned 50k checkpoints and
continue training to 150k under four basic channel controls plus one
matched-stat continuation control (\(N = 15\) seeds).

\begin{table}[t]
  \caption{Same-checkpoint continuation controls: cooperation when all modes start from identical learned 50k policies. At 50k, all values are identical (same checkpoint). By 150k, the mixed-motive branch ($f\!=\!3.5$) retains a clear learned-channel advantage, whereas the cooperation-dominant regime ($f\!=\!5.0$) shows little separation. The matched-stat \texttt{sender\_shuffle} row preserves per-sender token frequencies while breaking state-contingency. Mean $\pm$ SEM over \TIIInSeeds{} seeds; $\Delta$ column shows the paired learned$-$control gap at 150k with 95\% bootstrap CI.}
  \tablabel{sameckpt}
  \centering
  \small
  \begin{tabular}{lccccc}
    \toprule
    & \multicolumn{2}{c}{$f\!=\!3.5$} & \multicolumn{2}{c}{$f\!=\!5.0$} & \\
    \cmidrule(lr){2-3} \cmidrule(lr){4-5}
    Channel mode & 50k & 150k & 50k & 150k & $\Delta$ at 150k ($f\!=\!3.5$) \\
    \midrule
    \texttt{learned}        & $\TIIILearnedFk$ & $\mathbf{\TIIILearnedOfk}$ & $\TIIILearnedFkFV$ & $\TIIILearnedOfkFV$ & --- \\
    \texttt{indep\_random}  & $\TIIIIndepFk$ & $\TIIIIndepOfk$ & $\TIIIIndepFkFV$ & $\TIIIIndepOfkFV$ & $\TIIIIndepGapOfk\pp\;\TIIIIndepGapOfkCI$ \\
    \texttt{public\_random} & $\TIIIPubFk$ & $\TIIIPubOfk$ & $\TIIIPubFkFV$ & $\TIIIPubOfkFV$ & $\TIIIPubGapOfk\pp\;\TIIIPubGapOfkCI$ \\
    \texttt{always\_zero}   & $\TIIIZeroFk$ & $\TIIIZeroOfk$ & $\TIIIZeroFkFV$ & $\TIIIZeroOfkFV$ & $\TIIIZeroGapOfk\pp\;\TIIIZeroGapOfkCI$ \\
    \midrule
    \texttt{sender\_shuffle} & $\TIIILearnedFk$ & $\TSSshufCoopTFOfkPM$ & $\TIIILearnedFkFV$ & $\TSSshufCoopFVOfkPM$ & $\TSSgapTFOfk\pp\;\TSSgapTFOfkCI$ \\
    \bottomrule
  \end{tabular}
\end{table}

Because all branches start from the same 50k checkpoint, later
divergence isolates the effect of channel mode during continued training
from 50k to 150k. Operationally, this is a common warm-start
continuation protocol: each branch reloads the same 50k network weights
and resumes the learning-rate and entropy schedules from the same
episode count, but does not restore optimizer state. We therefore
interpret Table \ref{tab:sameckpt} as evidence about continued-training
channel-content effects under a shared continuation procedure rather
than as a complete measure of communication's value over the full
training run. By 150k at \(f\!=\!3.5\), learned messages retain paired
advantages of \(\TIIIZeroGapOfk\pp\) versus silence (95\% CI
\(\TIIIZeroGapOfkCI\), sign-flip \(p\TIIIZeroGapOfkP\)),
\(\TIIIPubGapOfk\pp\) versus public random (CI \(\TIIIPubGapOfkCI\),
\(p\TIIIPubGapOfkP\)), and \(\TIIIIndepGapOfk\pp\) versus independent
random (CI \(\TIIIIndepGapOfkCI\), \(p\TIIIIndepGapOfkP\)). All three
paired confidence intervals exclude zero. The basic zero/random control
modes all cluster tightly (standard deviations
\({\sim}0.04\)--\(0.08\)), while the learned condition shows
substantially higher variance (SD \({\sim}0.15\)), suggesting that
semantic content amplifies seed-level heterogeneity. At \(f\!=\!5.0\),
those basic four modes remain close together, while the matched-stat
\texttt{sender\_shuffle} branch later shows that reshuffling can even
outperform learned communication. The cleanest reading is therefore not
generic ``difficulty gating,'\,' but strategic value-of-information
gating: message content matters where agents face mixed incentives and a
noisy latent regime, and matters little where cooperation is already
individually rational.

Welfare tracks the same pattern, so the mixed-regime result is not just
a cooperation-only artifact.

\textbf{Caveat: channel-distribution shift.} These controls are not
perfectly semantic-clean. Replacing learned messages with zero or random
inputs changes the input distribution seen by a channel-conditioned
policy, so some of the learned-minus-control gap could reflect
adaptation cost under altered channel statistics rather than semantic
value alone. The magnitude of this confound is bounded by the fact that
all three control modes (which differ substantially in their input
statistics---silence, independent noise, and correlated noise) cluster
tightly at 150k, while only the learned branch diverges upward. We
therefore added a matched-stat continuation control that preserves each
sender's marginal token frequency while breaking state-contingency
(\texttt{sender\_shuffle}). At 150k and \(f\!=\!3.5\), learned still
exceeds sender-shuffle on average (53.6\% vs.~44.6\% cooperation;
\(+9.0\pp\) mean paired gap), but this matched-stat gap is heterogeneous
and its bootstrap CI overlaps zero. At \(f\!=\!5.0\), sender-shuffle
instead outperforms learned by \(15.9\pp\) on average, with a negative
bootstrap CI that excludes zero. The cleanest interpretation is
therefore intermediate: the learned-vs-zero/random advantage likely
combines genuine state-linked message value with broader structural
dependence on having a channel, and the mixed-regime semantic component
appears real but more modest under matched-stat controls.

Figure~\ref{fig-value-surface} makes the multiplier-wide version of this
point visible before we turn to the endpoint intervention suite.

\begin{figure}[t]

{\centering \includegraphics{main_files/figure-pdf/fig-value-surface-output-1.pdf}

}

\caption{\label{fig-value-surface}Full multiplier surface for three key
150k contrasts. Left: separately trained communication minus
no-communication endpoint gap. Middle: same-checkpoint learned minus
silence continuation gap. Right: frozen natural minus sender-shuffle
endpoint gap. The continued-learning content effect peaks in the upper
mixed regimes (\(f\!=\!2.5, 3.5\)), while the raw endpoint gap reverses
at \(f\!=\!5.0\) and sender-shuffling is most harmful to the natural
endpoint protocol at \(f\!=\!2.5, 3.5\).}

\end{figure}

This reveals a \textbf{strategic value-of-information} pattern:

\begin{itemize}[leftmargin=*,nosep]
  \item \textbf{Mixed-motive regime ($f\!=\!3.5$)}: Learned message content provides a clear advantage over noise and silence when private signals are noisy and unilateral cooperation is still costly.
  \item \textbf{Cooperation-dominant regime ($f\!=\!5.0$)}: Message content is largely irrelevant. Agents learn to cooperate regardless of channel mode.
\end{itemize}

Seed heterogeneity remains substantial even at 15 seeds, underscoring
the importance of multi-seed evaluation in MARL experiments. As shown in
Figure~\ref{fig-semantic-window}, the mixed-motive paired deltas are
mostly positive against all three controls, while the
cooperation-dominant raw outcomes overlap heavily across channel modes.

\begin{figure}[t]

{\centering \includegraphics{main_files/figure-pdf/fig-semantic-window-output-1.pdf}

}

\caption{\label{fig-semantic-window}Same-checkpoint continuation
controls at 150k episodes (15 seeds). \textbf{Left:} Paired
learned\(-\)control deltas at \(f\!=\!3.5\) for each seed.
\textbf{Right:} Raw per-seed cooperation at \(f\!=\!5.0\) across the
four channel modes. Black squares show mean \(\pm\) SEM.}

\end{figure}

\hypertarget{sec-frozen-endpoint}{%
\subsection{Frozen-checkpoint endpoint interventions are weak at
150k}\label{sec-frozen-endpoint}}

Same-checkpoint continuation shows that meaningful message content
affects the 50k-to-150k learning trajectory. To ask the narrower
endpoint question---how much the frozen learned policy still depends on
online messages---we evaluate the learned checkpoints under direct
message perturbations without further training.

\begin{table}[t]
  \caption[Frozen-checkpoint intervention effects at $f\!=\!3.5$.]{Frozen-checkpoint intervention effects in the mixed-motive regime ($f\!=\!3.5$). $\Delta P(C)$ in percentage points, computed as natural $-$ intervened: positive values mean natural messages helped, negative values mean the intervention \emph{outperformed} natural messages. Mean $\pm$ SEM over 15 seeds; 95\% bootstrap CI and sign consistency (fraction of seeds where natural $>$ intervention) shown for 150k.}
  \tablabel{frozen_endpoint}
  \centering
  \small
  \begin{tabular}{lcccl}
    \toprule
    Intervention & 50k & 150k & 95\% CI & Sign \\
    \midrule
    \texttt{zeros}          & $\TFRZerosFk \pm \TFRZerosFkSEM$ & $\TFRZerosOfk \pm \TFRZerosOfkSEM$ & $\TFRZerosOfkCI$ & $\TFRZerosOfkSign$ \\
    \texttt{fixed0}         & $\TFRFixZeroFk \pm \TFRFixZeroFkSEM$ & $\TFRFixZeroOfk \pm \TFRFixZeroOfkSEM$ & $\TFRFixZeroOfkCI$ & $\TFRFixZeroOfkSign$ \\
    \texttt{public\_random} & $\TFRPubFk \pm \TFRPubFkSEM$ & $\TFRPubOfk \pm \TFRPubOfkSEM$ & $\TFRPubOfkCI$ & $\TFRPubOfkSign$ \\
    \texttt{indep\_random}  & $\TFRIndepFk \pm \TFRIndepFkSEM$ & $\TFRIndepOfk \pm \TFRIndepOfkSEM$ & $\TFRIndepOfkCI$ & $\TFRIndepOfkSign$ \\
    \texttt{sender\_shuffle} & $\TFRShuffleFk \pm \TFRShuffleFkSEM$ & $\TFRShuffleOfk \pm \TFRShuffleOfkSEM$ & $\TFRShuffleOfkCI$ & $\TFRShuffleOfkSign$ \\
    \texttt{permute\_slots} & $\TFRPermuteFk \pm \TFRPermuteFkSEM$ & $\TFRPermuteOfk \pm \TFRPermuteOfkSEM$ & $\TFRPermuteOfkCI$ & $\TFRPermuteOfkSign$ \\
    \bottomrule
  \end{tabular}
\end{table}

Table \ref{tab:frozen_endpoint} reveals a mixed endpoint picture. At
50k, most interventions have near-zero effects: natural messages are not
cleanly better than replacements. By 150k, natural messages are on
average above several replacements, including silence (\texttt{zeros},
\(\TFRZerosOfk\pp\)) and independent random noise
(\texttt{indep\_random}, \(\TFRIndepOfk\pp\)), but these effects remain
modest and heterogeneous across seeds.

Critically, however, \texttt{sender\_shuffle} (\(\TFRShuffleOfk\pp\))
\emph{outperforms} natural messages at 150k---and does so consistently
across all 15 seeds (\(\TFRShuffleOfkSign\) seeds with natural \(>\)
shuffled; 95\% CI: \(\TFRShuffleOfkCI\pp\), sign-flip
\(p\TFRShuffleOfkP\)). Sender-shuffled messages preserve each sender's
marginal token distribution but break the link between message and
current state. That shuffled messages beat natural ones shows that the
frozen endpoint policy does not reliably benefit from the original
state-conditioned message stream. This does not imply that endpoint
messages are meaningless. Rather, it suggests that whatever endpoint
structure exists is brittle: it may be sender-indexed, poorly
calibrated, or otherwise misaligned with a robust shared code.
Sender-indexed conventions are one live explanation, but not the only
one: self-conditioning, temporal miscalibration, and more general
channel-conditioned overfitting remain consistent with this intervention
pattern. To test whether endpoint messages are still used causally at a
more local level, we next intervene on individual sender tokens while
holding the rest of the message vector fixed.

At the frozen 150k checkpoint, we therefore run a sender-specific causal
probe: for each stored greedy-evaluation state and true multiplier, we
hold raw observations and all other message slots fixed, force sender
\(i\)'s token to 0 or 1, and measure the induced change in receiver
\(j\)'s cooperation probability and greedy action.

\begin{table}[t]
  \caption{Frozen sender-specific causal probe at 150k. Non-self sender--receiver pairs only. For each true multiplier $f$, we hold the raw observation and remaining message slots fixed, force sender $i$'s token to 0 or 1, and average over seeds the resulting signed change in receiver cooperation probability, the absolute change, and the greedy-action flip rate relative to the natural message.}
  \tablabel{sender_causal}
  \centering
  \small
  \begin{tabular}{lcccc}
    \toprule
    $f$ & Signed $\Delta P(C)$ (pp) & Mean $|\Delta P(C)|$ (pp) & Flip rate force0 (\%) & Flip rate force1 (\%) \\
    \midrule
    $0.5$ & $\TSCDeltaZeroFive$ & $\TSCAbsZeroFive$ & $\TSCFlipZeroZeroFive$ & $\TSCFlipOneZeroFive$ \\
    $1.5$ & $\TSCDeltaOneFive$ & $\TSCAbsOneFive$ & $\TSCFlipZeroOneFive$ & $\TSCFlipOneOneFive$ \\
    $2.5$ & $\TSCDeltaTwoFive$ & $\TSCAbsTwoFive$ & $\TSCFlipZeroTwoFive$ & $\TSCFlipOneTwoFive$ \\
    $3.5$ & $\TSCDeltaThreeFive$ & $\TSCAbsThreeFive$ & $\TSCFlipZeroThreeFive$ & $\TSCFlipOneThreeFive$ \\
    $5.0$ & $\TSCDeltaFiveZero$ & $\TSCAbsFiveZero$ & $\TSCFlipZeroFiveZero$ & $\TSCFlipOneFiveZero$ \\
    \bottomrule
  \end{tabular}
\end{table}

Table \ref{tab:sender_causal} sharpens the endpoint interpretation. The
signed mean effect is near zero at every multiplier, including
\(\TSCDeltaThreeFive\pp\) at \(f\!=\!3.5\) and \(\TSCDeltaFiveZero\pp\)
at \(f\!=\!5.0\), but the mean absolute non-self sender--receiver effect
remains sizable: \(\TSCAbsTwoFive\pp\) at \(f\!=\!2.5\),
\(\TSCAbsThreeFive\pp\) at \(f\!=\!3.5\), and \(\TSCAbsFiveZero\pp\) at
\(f\!=\!5.0\). Forcing a single sender's token flips the receiver's
greedy action \(\TSCFlipZeroThreeFive\)--\(\TSCFlipOneThreeFive\)\% of
the time at \(f\!=\!3.5\) and
\(\TSCFlipZeroFiveZero\)--\(\TSCFlipOneFiveZero\)\% at \(f\!=\!5.0\).
Endpoint messages are therefore not inert; they exert heterogeneous
local causal effects that largely cancel in aggregate.

This endpoint use is not reducible to pure self-conditioning. In the
same 150k sender-specific probe, self pairs are somewhat stronger than
non-self pairs, but non-self effects remain substantial: at
\(f\!=\!3.5\), mean absolute \(|\Delta P(C)|\) is \(5.7\pp\) for self
pairs and \(5.0\pp\) for non-self pairs, with non-self greedy-action
flip rates still 12.2--12.9\%. At \(f\!=\!5.0\), the corresponding
values are \(4.0\pp\) versus \(3.8\pp\), with non-self flip rates
9.8--11.5\%. The message channel therefore appears to serve both
self-conditioning and genuine cross-agent influence at the endpoint.

These endpoint-causal effects are still much smaller than the
same-checkpoint continuation gaps (\(+12\)--\(19\pp\)). The synthesis is
therefore two-layered: meaningful communication shaped the learning
trajectory substantially, but the frozen policy uses messages weakly and
inconsistently online.

An exploratory mute-after-50k control remains consistent with structural
channel dependence but is only 3 seeds, so we leave it in
Appendix\textasciitilde E rather than treat it as a core result.

\hypertarget{synthesis-separating-continued-learning-from-endpoint-effects}{%
\subsection{Synthesis: separating continued-learning from endpoint
effects}\label{synthesis-separating-continued-learning-from-endpoint-effects}}

The dissociation between continued-learning value and endpoint protocol
coherence is not confined to the single \(f\!=\!3.5\) versus
\(f\!=\!5.0\) comparison. As Figure~\ref{fig-value-surface} already
previewed, the same-checkpoint continuation gap peaks in the upper mixed
regimes, while the frozen natural-versus-shuffle contrast is also
largest there.

Overall, the results show a dissociation. Same-checkpoint continuation
controls reveal continued-learning value, while frozen-endpoint
interventions ask a narrower question about online message use in a
fixed policy. In this EPGG case study:

\begin{enumerate}[leftmargin=*,nosep]
  \item \textbf{Continued-learning value survives in the mixed-motive regime, but weakens under matched-stat controls.} At $f\!=\!3.5$, learned messages outperform all control modes when forked from identical checkpoints (150k, Table \ref{tab:sameckpt}). By 150k, the learned branch remains $+19\pp$ above silence and $+12$--$14\pp$ above the random-channel controls. A matched-stat sender-shuffle branch still trails learned by $+9.0\pp$ on average at 150k, but this gap is appreciably smaller and more heterogeneous across seeds, indicating that the strongest learned-vs-control separation partly reflects broader channel dependence as well as state-linked message value.
  \item \textbf{Endpoint online use is much weaker and can be adverse.} The frozen-endpoint interventions are small relative to the same-checkpoint gaps, and \texttt{sender\_shuffle} improves performance. Yet the sender-specific causal probe still finds $\TSCAbsThreeFive\pp$ mean absolute non-self effects at $f\!=\!3.5$ despite a signed average near zero. Endpoint communication is therefore locally causal but not organized into a robust shared protocol. This maps directly to the \emph{aggregate cancellation} failure mode identified in the introduction: sender-indexed conventions with opposing polarities produce near-zero aggregate effects despite strong per-sender responses.
  \item \textbf{The $f\!=\!5.0$ reversal is consistent with training-trajectory differences rather than ongoing message-content effects.} In this EPGG setting, communication-trained agents cooperate $\pvGapFVOfkPP\pp$ less than no-communication agents at $f\!=\!5.0$ at 150k. Since same-checkpoint controls show no channel-mode separation at $f\!=\!5.0$, the reversal is more naturally attributed to training-trajectory differences and architectural dependence than to harmful endpoint semantics per se. This illustrates the \emph{strategic value-of-information gating} failure mode: evaluating only at $f\!=\!5.0$ (or averaging across regimes) would obscure the genuine content effect at $f\!=\!3.5$.
\end{enumerate}

The exploratory mute-after-50k control is additionally consistent with
structural dependence, but the same-checkpoint and frozen-endpoint
results already establish the central dissociation without relying on
this small-(N) add-on. Aggregate endpoint summaries can still miss
sender-conditioned structure because opposing per-sender effects cancel
in the mean. Appendix\textasciitilde C reports that observational
disaggregation and polarity analysis; the main identification rests on
the paired continuation controls and frozen causal probes above.

\hypertarget{sec-discussion}{%
\section{Discussion}\label{sec-discussion}}

\textbf{Summary.} Same-checkpoint continuation controls reveal
continued-learning value, while frozen-checkpoint interventions ask
whether the endpoint policy uses messages as a robust online code. In
this EPGG case study, learned messages retain a clear mixed-regime
advantage over silence and random channels during continued training,
but a matched-stat sender-shuffle control shrinks that advantage to
about \(+9\pp\), indicating a mixture of state-linked message value and
broader channel dependence. At the frozen endpoint, sender-shuffled
messages outperform natural ones while sender-specific probes still show
sizable local effects, so the endpoint code is not inert but brittle and
sender-indexed.

\textbf{Methodological point.} The main contribution is therefore
evaluative rather than game-specific: endpoint summaries alone can miss
a split between continued-learning value and endpoint protocol
coherence, especially under independent learners where non-stationarity
already complicates attribution
\citep{papoudakis2019dealing, oroojlooyjadid2022review}.

\textbf{Limitations.} The results are conditional on one game, a binary
vocabulary, four fixed-team agents with sender-labeled slots, and
independent PPO. The same-checkpoint design holds initialization fixed
but does not fully remove channel-distribution shift, and our
matched-stat sender-shuffle control still preserves only per-sender
marginals rather than richer temporal structure or a clean
self-versus-other decomposition. The continuation branches are also warm
starts rather than bitwise uninterrupted runs, so stitched online curves
are descriptive only; the main policy-quality evidence comes from greedy
checkpoint evaluation.

\textbf{Future work.} The most immediate next steps are stricter
matched-history controls, explicit self-versus-other message masking,
and multi-checkpoint branching to map when communication matters during
training rather than only from 50k to 150k.

\hypertarget{references}{%
\subsection*{References}\label{references}}
\addcontentsline{toc}{subsection}{References}

\renewcommand{\bibsection}{}
\bibliography{refs.bib}

\appendix

\clearpage

\hypertarget{sec-app-params}{%
\section{Appendix A: Full experimental
parameters}\label{sec-app-params}}

\begin{table}[H]
  \caption{Complete hyperparameter specification for all training runs.}
  \tablabel{hyperparams}
  \centering
  \small
  \begin{tabular}{ll}
    \toprule
    Parameter & Value \\
    \midrule
    \multicolumn{2}{l}{\textit{Environment}} \\
    Number of agents $N$ & 4 \\
    Episode length $T$ & 100 rounds \\
    Multiplier set $\mathcal{F}$ & $\{0.5, 1.5, 2.5, 3.5, 5.0\}$ \\
    Regime switch probability $\rho$ & 0.05 \\
    Signal noise $\sigma$ & 0.5 (per agent) \\
    Tremble probability $\varepsilon_{\text{tremble}}$ & 0.05 \\
    Endowment $c$ & 4 \\
    \midrule
    \multicolumn{2}{l}{\textit{Architecture}} \\
    Policy/value networks & Separate 2-layer MLPs (action, message, value) \\
    Hidden size & 64 \\
    Recurrent state & None \\
    Observation wrapper & $[\fhat,\ \text{endowment},\ \text{last coop. frac.},\ \text{own last action},\ \text{EWMA coop.},\ \text{received one-hot messages}]$ \\
    Vocabulary size $|M|$ & 2 (binary) \\
    Parameters per network & ${\sim}5$K scale \\
    \midrule
    \multicolumn{2}{l}{\textit{PPO}} \\
    Discount $\gamma$ & 0.99 \\
    GAE $\lambda$ & 0.95 \\
    Clip ratio & 0.2 \\
    Reward scale & 20.0 \\
    \midrule
    \multicolumn{2}{l}{\textit{Schedules}} \\
    Learning rate & $3 \times 10^{-4} \to 10^{-5}$ (cosine) \\
    Action entropy coeff. & $0.01 \to 0.001$ (linear) \\
    Message entropy coeff. & $0.01 \to 0.0$ (linear) \\
    \midrule
    \multicolumn{2}{l}{\textit{Training}} \\
    Total episodes & 150,000 \\
    Seeds & 15 independent seeds \\
    Checkpoint interval & 50,000 episodes \\
    \midrule
    \multicolumn{2}{l}{\textit{Evaluation}} \\
    Action selection & Greedy (argmax) \\
    Evaluation episodes & 300 per (checkpoint, seed, condition) \\
    Evaluation seed & 9001 (fixed) \\
    \bottomrule
  \end{tabular}
\end{table}

\hypertarget{sec-app-per-seed}{%
\section{Appendix B: Per-seed breakdown of main
results}\label{sec-app-per-seed}}

\begin{table}[H]
  \caption{Summary statistics for cooperation rates and communication gap at $f\!=\!3.5$ across training ($n = \pvPSnSeedsFk$ seeds).}
  \tablabel{per_seed_summary}
  \centering
  \small
  \begin{tabular}{llccc}
    \toprule
    & & \multicolumn{3}{c}{Checkpoint (episodes)} \\
    \cmidrule(lr){3-5}
    Metric & Statistic & 50k & 100k & 150k \\
    \midrule
    \multirow{3}{*}{Comm $P(C)$}
      & Mean & $\pvPScommMeanFk$ & $\pvPScommMeanHk$ & $\pvPScommMeanOfk$ \\
      & Min & $\pvPScommMinFk$ & $\pvPScommMinHk$ & $\pvPScommMinOfk$ \\
      & Max & $\pvPScommMaxFk$ & $\pvPScommMaxHk$ & $\pvPScommMaxOfk$ \\
    \midrule
    \multirow{3}{*}{Comm gap $\Delta$}
      & Mean & $\pvPSgapMeanFk$ & $\pvPSgapMeanHk$ & $\pvPSgapMeanOfk$ \\
      & Min & $\pvPSgapMinFk$ & $\pvPSgapMinHk$ & $\pvPSgapMinOfk$ \\
      & Max & $\pvPSgapMaxFk$ & $\pvPSgapMaxHk$ & $\pvPSgapMaxOfk$ \\
    \bottomrule
  \end{tabular}
\end{table}

\hypertarget{sec-app-alignment}{%
\section{Appendix C: Aggregate endpoint metrics and polarity
alignment}\label{sec-app-alignment}}

The frozen sender-specific interventions in the main text show that
endpoint messages can still have heterogeneous local causal effects.
This appendix adds the descriptive complement: aggregate
receiver-semantics summaries can nevertheless look near zero once they
average across senders.

\textbf{Important caveat.} The receiver-semantics effects reported here
are \emph{observational} (conditioned on sender identity), not
interventional. They describe sender-conditioned structure but do not
isolate a causal mechanism.

\hypertarget{per-sender-disaggregation-reveals-hidden-structure}{%
\subsection{Per-sender disaggregation reveals hidden
structure}\label{per-sender-disaggregation-reveals-hidden-structure}}

For each receiver \(j\) and sender \(i\), we compute the per-sender
token effect within observation bins:
\begin{equation}\protect\hypertarget{eq-per-sender}{}{
  \delta_{i \to j}(\fhat_{\mathrm{bin}}) = \pcoopgiven{m_i\!=\!1, \fhat_{\mathrm{bin}}} - \pcoopgiven{m_i\!=\!0, \fhat_{\mathrm{bin}}}
}\label{eq-per-sender}\end{equation} and compare it to the aggregate
effect, which pools across all senders:
\begin{equation}\protect\hypertarget{eq-agg}{}{
  \delta_{\text{agg}}(\fhat_{\mathrm{bin}}) = P(C \mid \exists\, k: m_k\!=\!1,\; \fhat_{\mathrm{bin}}) - P(C \mid \forall\, k: m_k\!=\!0,\; \fhat_{\mathrm{bin}})
}\label{eq-agg}\end{equation}

\begin{table}[H]
  \caption{Aggregate versus per-sender token effects across training. Aggregate effects are near zero due to convention cancellation, while sender-conditioned effects (mean $|\delta_{i \to j}|$) remain large. Mean over 15 seeds.}
  \tablabel{disagg}
  \centering
  \small
  \begin{tabular}{lccc}
    \toprule
    & \multicolumn{3}{c}{Checkpoint (episodes)} \\
    \cmidrule(lr){2-4}
    Metric & 50k & 100k & 150k \\
    \midrule
    Aggregate token effect     & $\TVaggFk$ & $\TVaggHk$ & $\TVaggOfk$ \\
    Per-sender token effect    & $\TVpsFk$ & $\TVpsHk$  & $\TVpsOfk$ \\
    Ratio (per-sender / |agg.|)  & $\TVratioFk$ & $\TVratioHk$ & $\TVratioOfk$ \\
    \bottomrule
  \end{tabular}
\end{table}

At 50k, the aggregate token effect is only \(\pvAggFkPP\pp\), while the
per-sender effect is \(\pvPsFkPP\pp\). The same discrepancy persists at
150k: sender-conditioned effects remain large even though the aggregate
summary stays near zero.

\begin{figure}[t]

{\centering \includegraphics{main_files/figure-pdf/fig-disagg-output-1.pdf}

}

\caption{\label{fig-disagg}Aggregate vs.~per-sender token effects across
training. The aggregate effect (signed mean \(\delta\)) is near zero at
all checkpoints due to convention fragmentation. The per-sender effect
(mean \(|\delta_{i \to j}|\)) reveals strong sender-conditioned
responses that aggregate metrics hide.}

\end{figure}

\hypertarget{polarity-alignment-patterns}{%
\subsection{Polarity alignment
patterns}\label{polarity-alignment-patterns}}

The aggregate cancellation occurs because sender-indexed conventions
often adopt opposing polarities. We summarize polarity by the sign of
\(P(m\!=\!1 \mid \fhat \ge 4.5) - P(m\!=\!1 \mid \fhat < 1.5)\).

Across \TVInSeeds{} seeds, alignment trajectories are heterogeneous:
some seeds show stable alignment, while others flip or lose it between
50k and 150k. This confirms that convention formation in
independent-learner MARL is stochastic rather than converging reliably
to one shared code.

\hypertarget{sec-app-separate}{%
\section{Appendix D: Separate-training channel controls (exploratory, 3
seeds)}\label{sec-app-separate}}

This exploratory experiment trains agents from scratch under four
channel modes for 50k episodes: \texttt{learned} (sender's actual
message), \texttt{always\_zero} (all sender slots forced to 0),
\texttt{indep\_random} (sender slots randomized independently each
round), and \texttt{public\_random} (all sender slots share the same
random bit each round). The learned condition uses 5 seeds; controls use
3.

\begin{table}[H]
  \caption{Cooperation at $f\!=\!3.5$ after 50k episodes under different channel modes, trained from scratch. Exploratory (3 seeds per control, 5 for learned).}
  \tablabel{separate_controls}
  \centering
  \small
  \begin{tabular}{lc}
    \toprule
    Channel mode & $\pcoopgiven{f\!=\!3.5}$ \\
    \midrule
    \texttt{learned}        & $\mathbf{\TIILearned}$ \\
    \texttt{indep\_random}  & $\TIIIndep$ \\
    \texttt{always\_zero}   & $\TIIZero$ \\
    \texttt{public\_random} & $\TIIPub$ \\
    \bottomrule
  \end{tabular}
\end{table}

Because this design confounds channel mode with the entire training
trajectory, it should be read as suggestive rather than decisive. The
table is provided for transparency.

\hypertarget{sec-app-mute}{%
\section{Appendix E: Mute experiment details (exploratory, 3
seeds)}\label{sec-app-mute}}

The mute-after-50k experiment trains with learned communication until
50k, then delivers silence (all-zero messages) for the remaining 100k
episodes while senders continue to compute and update message outputs.
This tests whether the action policy can function without the message
inputs it was trained on. With only 3 seeds, this result is exploratory.

\begin{table}[H]
  \caption{Cooperation and welfare at $f\!=\!3.5$ at 150k. Muting after 50k produces cooperation \emph{below} the no-communication baseline. Exploratory (3 seeds for mute condition; 15 for comm and no-comm).}
  \tablabel{mute}
  \centering
  \small
  \begin{tabular}{lcc}
    \toprule
    Condition & $\pcoopgiven{f\!=\!3.5}$ & Welfare ($f\!=\!3.5$) \\
    \midrule
    \texttt{comm\_symm} (channel stays on)    & $\TIVcommCoop$ & $\TIVcommWelf$ \\
    \texttt{no\_comm\_symm} (never had channel) & $\TIVncCoop$ & $\TIVncWelf$ \\
    \texttt{mute-after-50k}                   & $\TIVmuteCoop$ & $\TIVmuteWelf$ \\
    \bottomrule
  \end{tabular}
\end{table}

The cooperation deficit (\(-\pvMuteBelowPP\pp\) relative to
no-communication) is consistent with structural dependence, but with
only 3 seeds it should be treated as suggestive rather than definitive.

\hypertarget{sec-app-online}{%
\section{Appendix F: Online training curves}\label{sec-app-online}}

\begin{figure}[t]

{\centering \includegraphics{../../outputs/eval/paper_strengthen/iter9_online_vs_checkpoint_15seeds/online_vs_checkpoint_f35_f50.png}

}

\caption{\label{fig-coop-trajectory}Dense online training curves with
sparse checkpoint anchors for the two most informative multipliers.
Solid lines show sampled training-window behavior from 1k to 150k;
shaded bands are SEM across 15 seeds. Dashed lines with dots show greedy
checkpoint evaluation at 25k, 50k, 100k, and 150k. The vertical marker
at 50k indicates the continuation boundary between the initial 1k--50k
run and the 55k--150k extension.}

\end{figure}

These curves mix sampled training-window behavior with greedy checkpoint
evaluation, so they should be read descriptively rather than as the
paper's main identification result. The sharp break at 50k marks the
continuation seam between the original 1k--50k run and the resumed
55k--150k extension.

\newpage
\section*{NeurIPS Paper Checklist}

\begin{enumerate}

\item {\bf Claims}
    \item[] Answer: \answerYes{}
    \item[] Justification: The abstract and introduction clearly state four contributions supported by the experimental results in the main body.

\item {\bf Limitations}
    \item[] Answer: \answerYes{}
    \item[] Justification: The \emph{Discussion} section discusses limitations including single game, minimal vocabulary, PPO coupling, and seed heterogeneity.

\item {\bf Theory assumptions and proofs}
    \item[] Answer: \answerNA{}
    \item[] Justification: This is an empirical paper; no theoretical results are presented.

\item {\bf Experimental result reproducibility}
    \item[] Answer: \answerYes{}
    \item[] Justification: Appendix A provides complete hyperparameters. Code will be released upon acceptance.

\item {\bf Open access to data and code}
    \item[] Answer: \answerYes{}
    \item[] Justification: Anonymized code and evaluation scripts will be provided as supplementary material. Full code release upon acceptance.

\item {\bf Experimental setting/details}
    \item[] Answer: \answerYes{}
    \item[] Justification: The \emph{Setup} section and Appendix A provide all training and evaluation details.

\item {\bf Experiment statistical significance}
    \item[] Answer: \answerYes{}
    \item[] Justification: We report mean $\pm$ SEM across 15 seeds, bootstrap confidence intervals (10,000 resamples) for effect sizes, and exact paired sign-flip tests for primary paired comparisons.

\item {\bf Experiments compute resources}
    \item[] Answer: \answerYes{}
    \item[] Justification: Experiments ran on an Apple M1 MacBook Pro (10 cores, 16GB RAM) and an IWR cluster node (quadxeon7). Each 150k-episode training run takes approximately 9 hours. Total compute for all experiments (including pilots, ablations, and the 15-seed expansion) is approximately 1500 CPU-hours.

\item {\bf Code of ethics}
    \item[] Answer: \answerYes{}
    \item[] Justification: This research studies cooperation in abstract game-theoretic settings with no human subjects, sensitive data, or dual-use concerns.

\item {\bf Broader impacts}
    \item[] Answer: \answerNA{}
    \item[] Justification: This is basic research on emergent communication in abstract multi-agent systems. We do not foresee direct negative societal impacts.

\item {\bf Safeguards}
    \item[] Answer: \answerNA{}
    \item[] Justification: No models or datasets with misuse risk are released.

\item {\bf Licenses for existing assets}
    \item[] Answer: \answerNA{}
    \item[] Justification: No pre-existing datasets or models are used. The codebase builds on standard PyTorch.

\item {\bf New assets}
    \item[] Answer: \answerYes{}
    \item[] Justification: We will release the complete codebase (training, evaluation, analysis) and trained checkpoints upon acceptance.

\item {\bf Crowdsourcing and research with human subjects}
    \item[] Answer: \answerNA{}
    \item[] Justification: No human subjects involved.

\item {\bf Institutional review board (IRB) approvals or equivalent for research with human subjects}
    \item[] Answer: \answerNA{}
    \item[] Justification: No human subjects involved.

\item {\bf Declaration of LLM usage}
    \item[] Answer: \answerNA{}
    \item[] Justification: No LLMs are used as part of the core research methodology. LLMs were used only for writing assistance.

\end{enumerate}




\end{document}
